% Supplementary manuscript note: exact initialization of P0 and information sufficiency.
% This file is publication material and may be folded into fatigue_probability_main.tex
% during the next manuscript consolidation.

\section*{Supplementary note: initial-state generation of the probability measure}

For a fixed-left finite chain, define the ordered initial spacings and spacing rates by
\begin{equation}
\lambda_i(0)=x_i(0)-x_{i-1}(0),
\qquad
c_i(0)=\dot x_i(0)-\dot x_{i-1}(0).
\end{equation}
Because \(x_0=\dot x_0=0\), the ordered arrays \(\{\lambda_i(0),c_i(0)\}_{i=1}^M\) are equivalent to the complete deterministic microscopic state \(\Gamma_0\). The initial empirical phase-space and spacing measures are therefore generated exactly as
\begin{equation}
F_0(\lambda,c)
=\frac1M\sum_{i=1}^{M}
\delta[\lambda-\lambda_i(0)]\delta[c-c_i(0)],
\end{equation}
\begin{equation}
P_{\lambda,0}(\lambda)
=\frac1M\sum_{i=1}^{M}
\delta[\lambda-\lambda_i(0)].
\end{equation}
For an initial full-state measure \(\mu_0\), the corresponding ensemble forms follow by integrating these push-forwards over \(\mu_0(d\Gamma_0)\). Thus no named probability family is required at the initial time.

For a perfectly homogeneous static state under initial normalized force \(q_0=\sigma_0/E\), the stable equilibrium stretch satisfies
\begin{equation}
\phi'(\lambda_{\rm eq})=q_0,
\qquad
\phi''(\lambda_{\rm eq})>0.
\end{equation}
With zero initial spacing rates,
\begin{equation}
P_{\lambda,0}(\lambda)=\delta(\lambda-\lambda_{\rm eq}),
\qquad
F_0(\lambda,c)=\delta(\lambda-\lambda_{\rm eq})\delta(c).
\end{equation}
The delta form is therefore a consequence of the specified homogeneous deterministic initial condition rather than a statistical ansatz.

A one-point marginal is not, however, an exact autonomous initialization in the general chain. The spacing acceleration contains neighbor information,
\begin{equation}
\ddot\lambda_i
=\phi'(\lambda_{i+1})-2\phi'(\lambda_i)+\phi'(\lambda_{i-1}),
\end{equation}
so a permutation of the same set of spacing/rate pairs preserves both \(P_0\) and the one-point \(F_0(\lambda,c)\) while changing the deterministic accelerations. For a smooth test function \(g\),
\begin{equation}
\frac{d^2}{d\tau^2}
\left[\frac1M\sum_i g(\lambda_i)\right]
=
\frac1M\sum_i
\left[g''(\lambda_i)c_i^2+g'(\lambda_i)\ddot\lambda_i\right].
\end{equation}
Hence states with identical one-point \(F_0\) can have different second time derivatives of the same probability moment.

A direct generalized-LJ counterexample using \(m=12.19\), \(n=6\), zero initial rates, and two spatial permutations of the same eight spacings gives identical one-point \(P_0\) and \(F_0\) but an empirical spacing-distribution KS distance of \(0.375\) at normalized time \(\tau=1\). For \(g(\lambda)=\lambda^2\), the two initial second derivatives are approximately \(-3.51299\times10^{-3}\) and \(-1.85653\times10^{-3}\), respectively.

The exact finite-chain initialization hierarchy is therefore
\begin{equation}
\text{ordered }\{\lambda_i(0),c_i(0)\}
\Longleftrightarrow \Gamma_0
\Longrightarrow F_0(\lambda,c)
\Longrightarrow P_0(\lambda),
\end{equation}
with no general reverse implication. The full labeled initial state, or a probability measure over that state, remains the exact reference. Reduced initial probability data become sufficient only after an additional assumption removes or reconstructs the discarded spatial correlations.
